\chapter{Implémentation des modules administratifs et analytiques}

\section*{Introduction}
\begin{chapintro}
Ce chapitre décrit la réalisation du sixième et dernier sprint, qui élève la
plateforme au niveau enterprise sur deux axes complémentaires : d'un côté, un
\textbf{système RBAC (Role-Based Access Control) entièrement personnalisable}
permettant à l'administrateur de créer des rôles sur mesure et d'affecter des
permissions atomiques granulaires ; de l'autre, un \textbf{journal d'audit
enterprise} traçant chaque action avec ses valeurs avant/après, sa sévérité
automatique, l'adresse IP et le rôle de l'acteur, présenté via un tableau de
bord analytique interactif avec export PDF.
\end{chapintro}

\section{Sprint 6 : Sécurité enterprise et journal d'audit analytique}

% ============================================================
\subsection{Sprint Backlog}

\begin{table}[H]
  \centering
  \caption{Sprint Backlog — Sprint 6}
  \label{tab:backlog_s6}
  \begin{tabularx}{\textwidth}{C{1.2cm}L{5.8cm}L{2.2cm}C{2cm}C{1.5cm}}
    \toprule
    \thead{ID} & \thead{User Story} & \thead{Acteur} & \thead{Priorité} & \thead{Points} \\
    \midrule
    \rowcolor{lightgray}
    US-29 & Créer et gérer des rôles personnalisés        & Admin        & Must   &  8 \\
    US-30 & Affecter des permissions atomiques aux rôles   & Admin        & Must   &  8 \\
    \rowcolor{lightgray}
    US-31 & Consulter la matrice globale des permissions   & Admin        & Should &  5 \\
    US-32 & Consulter le journal d'audit enrichi           & Admin/Member & Must   &  8 \\
    \rowcolor{lightgray}
    US-33 & Analyser l'activité via le dashboard analytics & Admin        & Should &  8 \\
    US-34 & Exporter le journal d'audit en PDF             & Admin        & Could  &  5 \\
    \midrule
    \multicolumn{4}{r}{\textbf{Total :}} & \textbf{42} \\
    \bottomrule
  \end{tabularx}
\end{table}

% ============================================================
\subsection{Conception}

\subsubsection{Diagramme de cas d'utilisation}

\begin{figure}[H]
  \centering
  \IfFileExists{figures/sprint6_use_case.png}{%
    \includegraphics[width=0.88\textwidth]{figures/sprint6_use_case.png}%
  }{%
    \fbox{\parbox{0.88\textwidth}{\centering\vspace{1.8cm}
      \textbf{[IMAGE A INSÉRER]}\\[0.4em]
      \texttt{figures/sprint6\_use\_case.png}\\[0.3em]
      \textit{Cas d'utilisation — RBAC + Journal d'audit : Admin, Member, Système}
    \vspace{1.8cm}}}%
  }
  \caption{Diagramme de cas d'utilisation — Sprint 6}
  \label{fig:uc_s6}
\end{figure}

\subsubsection{Diagramme de séquence — Journal d'audit temps réel}

\begin{figure}[H]
  \centering
  \IfFileExists{figures/sprint6_sequence_audit.png}{%
    \includegraphics[width=0.92\textwidth]{figures/sprint6_sequence_audit.png}%
  }{%
    \fbox{\parbox{0.92\textwidth}{\centering\vspace{2cm}
      \textbf{[IMAGE A INSÉRER]}\\[0.4em]
      \texttt{figures/sprint6\_sequence\_audit.png}\\[0.3em]
      \textit{Séquence : Action utilisateur → Controller → ActivityLogService →
      activity\_logs (old/new values, IP, severity) → broadcast WebSocket →
      AuditLogPage}
    \vspace{2cm}}}%
  }
  \caption{Diagramme de séquence — Pipeline journal d'audit enterprise}
  \label{fig:seq_audit}
\end{figure}

\subsubsection{Diagramme de classes}

\begin{figure}[H]
  \centering
  \IfFileExists{figures/sprint6_class.png}{%
    \includegraphics[width=0.9\textwidth]{figures/sprint6_class.png}%
  }{%
    \fbox{\parbox{0.9\textwidth}{\centering\vspace{1.8cm}
      \textbf{[IMAGE A INSÉRER]}\\[0.4em]
      \texttt{figures/sprint6\_class.png}\\[0.3em]
      \textit{Classes : ActivityLog, ActivityLogService, ActivityFeedController,
      ActivityAnalyticsController, Role, Permission}
    \vspace{1.8cm}}}%
  }
  \caption{Diagramme de classes — Sprint 6}
  \label{fig:class_s6}
\end{figure}

\subsubsection{Architecture RBAC — Rôles et permissions}

Le système RBAC de TAC-TIC IntelliPM opère sur \textbf{deux niveaux distincts} :

\begin{table}[H]
  \centering
  \caption{Architecture RBAC à deux niveaux}
  \label{tab:rbac_arch}
  \begin{tabularx}{\textwidth}{L{3.5cm}L{3cm}X}
    \toprule
    \thead{Niveau} & \thead{Portée} & \thead{Mécanisme} \\
    \midrule
    \rowcolor{lightgray}
    Rôle système & Application & Spatie Permission — rôles \texttt{admin},
    \texttt{member}, \texttt{client} gérés dans la table \texttt{roles} \\
    Permission atomique & Fonctionnalité & 35 permissions granulaires (ex.
    \texttt{task.create}, \texttt{project.delete}, \texttt{role.view})
    pouvant être attribuées à n'importe quel rôle \\
    \rowcolor{lightgray}
    Rôle projet & Projet & Pivot \texttt{project\_user} avec colonne \texttt{role}
    (\texttt{owner}, \texttt{manager}, \texttt{developer}, \texttt{viewer}) \\
    \bottomrule
  \end{tabularx}
\end{table}

\subsubsection{Schéma enrichi de la table \texttt{activity\_logs}}

Huit nouvelles colonnes ont été ajoutées à la table existante pour constituer
un audit trail conforme aux exigences enterprise :

\begin{lstlisting}[style=sqlstyle,caption={Colonnes ajoutées — migration Sprint 6},
  label={lst:migration_s6}]
-- Migration : 2026_05_22_100000_add_enterprise_fields_to_activity_logs.php
ALTER TABLE activity_logs ADD COLUMN old_values    JSON         NULL;
ALTER TABLE activity_logs ADD COLUMN new_values    JSON         NULL;
ALTER TABLE activity_logs ADD COLUMN changed_fields JSON        NULL;
ALTER TABLE activity_logs ADD COLUMN ip_address    VARCHAR(45)  NULL;
ALTER TABLE activity_logs ADD COLUMN user_agent    VARCHAR(500) NULL;
ALTER TABLE activity_logs ADD COLUMN severity      VARCHAR(20)  DEFAULT 'info';
ALTER TABLE activity_logs ADD COLUMN source        VARCHAR(50)  DEFAULT 'web';
ALTER TABLE activity_logs ADD COLUMN actor_role    VARCHAR(100) NULL;

-- Index de performance :
CREATE INDEX idx_activity_logs_action            ON activity_logs(action);
CREATE INDEX idx_activity_logs_created_at        ON activity_logs(created_at);
CREATE INDEX idx_activity_logs_project_sev_date  ON activity_logs(project_id, severity, created_at);
\end{lstlisting}

\subsubsection{Détection automatique de sévérité}

\begin{table}[H]
  \centering
  \caption{Règles de sévérité automatique — ActivityLogService}
  \label{tab:severity_rules}
  \begin{tabularx}{\textwidth}{L{2.5cm}L{4cm}X}
    \toprule
    \thead{Sévérité} & \thead{Actions déclenchantes} & \thead{Utilisation UI} \\
    \midrule
    \rowcolor{lightgray}
    \textbf{critical} & \texttt{task.deleted}, \texttt{project.deleted} &
      Bande latérale rouge, badge rouge dans la timeline \\
    \textbf{warning}  & \texttt{task.assigned}, \texttt{project.member\_*},
      \texttt{sprint.*} & Bande latérale ambrée \\
    \rowcolor{lightgray}
    \textbf{info}     & Toutes les autres actions &
      Bande latérale bleue (défaut) \\
    \bottomrule
  \end{tabularx}
\end{table}

\subsubsection{Analytics d'activité — Algorithmes de détection d'anomalies}

Le contrôleur \texttt{ActivityAnalyticsController} calcule cinq métriques avancées
mises en cache pendant 5 minutes :

\begin{table}[H]
  \centering
  \caption{Métriques du dashboard analytics — calcul et seuils}
  \label{tab:analytics_metrics}
  \begin{tabularx}{\textwidth}{L{3.5cm}L{3cm}X}
    \toprule
    \thead{Métrique} & \thead{Seuil} & \thead{Algorithme SQL} \\
    \midrule
    \rowcolor{lightgray}
    Tâches instables & \textgreater{}5 mises à jour en 7\,j &
      \texttt{GROUP BY entity\_id HAVING count(*) > 5} \\
    Membres inactifs & 0 activité en 14\,j &
      \texttt{LEFT JOIN} sur \texttt{project\_members} excluant les acteurs récents \\
    \rowcolor{lightgray}
    Pics d'activité  & Count \textgreater{} 2× moyenne 28\,j &
      Comparaison par collection Eloquent après agrégation journalière \\
    Score collaboration & \% actions multi-utilisateurs × 2,5 (max 100) &
      \texttt{WHERE action IN ('task.assigned', 'comment\_added', ...)} \\
    \rowcolor{lightgray}
    Tendance 28\,j   & Série temporelle complète &
      \texttt{DATE(created\_at)}, jours manquants remplis à 0 \\
    \bottomrule
  \end{tabularx}
\end{table}

% ============================================================
\subsection{Réalisation}

\subsubsection{Page Rôles et permissions}

\begin{figure}[H]
  \centering
  \IfFileExists{figures/sprint6_01_roles.png}{%
    \includegraphics[width=0.88\textwidth]{figures/sprint6_01_roles.png}%
  }{%
    \fbox{\parbox{0.88\textwidth}{\centering\vspace{2cm}
      \textbf{[IMAGE A INSÉRER]}\\[0.4em]
      \texttt{figures/sprint6\_01\_roles.png}\\[0.3em]
      \textit{Page Rôles : liste des rôles avec nombre d'utilisateurs et de permissions,
      bouton créer rôle, badges colorés}
    \vspace{2cm}}}%
  }
  \caption{Page de gestion des rôles — Vue administrateur}
  \label{fig:ui_roles}
\end{figure}

\subsubsection{Matrice des permissions}

\begin{figure}[H]
  \centering
  \IfFileExists{figures/sprint6_02_permissions.png}{%
    \includegraphics[width=0.88\textwidth]{figures/sprint6_02_permissions.png}%
  }{%
    \fbox{\parbox{0.88\textwidth}{\centering\vspace{2cm}
      \textbf{[IMAGE A INSÉRER]}\\[0.4em]
      \texttt{figures/sprint6\_02\_permissions.png}\\[0.3em]
      \textit{Matrice : lignes = permissions atomiques, colonnes = rôles,
      cases = toggle actif/inactif avec sauvegarde instantanée}
    \vspace{2cm}}}%
  }
  \caption{Matrice globale des permissions}
  \label{fig:ui_permissions}
\end{figure}

\subsubsection{Journal d'audit enterprise — Onglet Timeline}

\begin{figure}[H]
  \centering
  \IfFileExists{figures/sprint6_03_audit_timeline.png}{%
    \includegraphics[width=0.88\textwidth]{figures/sprint6_03_audit_timeline.png}%
  }{%
    \fbox{\parbox{0.88\textwidth}{\centering\vspace{2cm}
      \textbf{[IMAGE A INSÉRER]}\\[0.4em]
      \texttt{figures/sprint6\_03\_audit\_timeline.png}\\[0.3em]
      \textit{Timeline groupée par date : bande sévérité colorée, humanisation naturelle,
      panneau avant/après expandable, badge WebSocket temps réel}
    \vspace{2cm}}}%
  }
  \caption{Journal d'audit — Timeline enterprise avec valeurs avant/après}
  \label{fig:ui_audit_timeline}
\end{figure}

L'humanisation des événements est assurée par la bibliothèque
\texttt{projectActivityHumanize.js} qui transforme les actions brutes en phrases
naturelles françaises :

\begin{lstlisting}[style=jsstyle,caption={Humanisation naturelle des actions d'audit},
  label={lst:humanize}]
// Entrée brute :
{ action: "task.updated", meta: { task_title: "Auth JWT",
  changed: ["priority"] }, old: { priority: "medium" },
  new: { priority: "high" } }

// Sortie humanisée :
"Fatma a modifié la priorité de « Auth JWT »"
// Panneau avant/après :
//   Priorité :  Moyenne  →  Haute
\end{lstlisting}

\subsubsection{Journal d'audit enterprise — Onglet Analytics}

\begin{figure}[H]
  \centering
  \IfFileExists{figures/sprint6_04_audit_analytics.png}{%
    \includegraphics[width=0.88\textwidth]{figures/sprint6_04_audit_analytics.png}%
  }{%
    \fbox{\parbox{0.88\textwidth}{\centering\vspace{2.2cm}
      \textbf{[IMAGE A INSÉRER]}\\[0.4em]
      \texttt{figures/sprint6\_04\_audit\_analytics.png}\\[0.3em]
      \textit{Dashboard analytics : 5 KPI cards, Sparkline 28j, Top 5 utilisateurs HBar,
      DonutStat sévérité, alertes anomalies (tâches instables, membres inactifs, pics)}
    \vspace{2.2cm}}}%
  }
  \caption{Tableau de bord analytics du journal d'audit}
  \label{fig:ui_audit_analytics}
\end{figure}

\begin{table}[H]
  \centering
  \caption{Widgets du dashboard analytics — implémentation CSS/SVG}
  \label{tab:analytics_widgets}
  \begin{tabularx}{\textwidth}{L{3.5cm}L{3cm}X}
    \toprule
    \thead{Widget} & \thead{Technologie} & \thead{Données affichées} \\
    \midrule
    \rowcolor{lightgray}
    5 cartes KPI   & HTML/CSS           & Activités 30\,j/7\,j, critiques,
                                          utilisateurs, collaboration \\
    Sparkline 28\,j & SVG \texttt{polyline} + dégradé & Tendance journalière
                                          (jours vides = 0) \\
    \rowcolor{lightgray}
    Top utilisateurs & Barres \texttt{div} (HBar) & 5 acteurs les plus actifs (30\,j)
                                          avec avatar \\
    Distribution sévérité & CSS \texttt{conic-gradient} & Répartition info/warning/critical \\
    \rowcolor{lightgray}
    Alertes anomalies & Liste colorée    & Tâches instables, membres inactifs,
                                          pics d'activité \\
    \bottomrule
  \end{tabularx}
\end{table}

\noindent\textit{Remarque :} tous les graphiques sont implémentés en CSS/SVG pur,
sans dépendance externe (pas de Recharts, Chart.js ou D3). Cela réduit la taille
du bundle de \textasciitilde{}150\,ko.

\subsubsection{Rapport PDF — Page d'impression dédiée}

\begin{figure}[H]
  \centering
  \IfFileExists{figures/sprint6_05_audit_pdf.png}{%
    \includegraphics[width=0.78\textwidth]{figures/sprint6_05_audit_pdf.png}%
  }{%
    \fbox{\parbox{0.78\textwidth}{\centering\vspace{1.8cm}
      \textbf{[IMAGE A INSÉRER]}\\[0.4em]
      \texttt{figures/sprint6\_05\_audit\_pdf.png}\\[0.3em]
      \textit{Page PDF A4 : en-tête TAC-TIC, synthèse analytique, tableau complet
      des événements avec sévérité et module colorés}
    \vspace{1.8cm}}}%
  }
  \caption{Rapport PDF du journal d'audit}
  \label{fig:ui_audit_pdf}
\end{figure}

L'export PDF est implémenté via une page React standalone (\texttt{AuditReportPage.jsx})
accessible à \texttt{/audit-log/report?...params}. La page charge les données
depuis l'API, rend le rapport, puis déclenche \texttt{window.print()} automatiquement.
Le format A4 est configuré avec \texttt{@media print} et \texttt{@page \{ size: A4 \}}.

\subsubsection{Sidebar SaaS redesign}

La barre de navigation latérale a été entièrement redessinée dans un style
inspiré des plateformes \textbf{Linear} et \textbf{Notion}, avec trois innovations
majeures par rapport à la version précédente :

\begin{table}[H]
  \centering
  \caption{Comparaison Sidebar v1 vs v2 (SaaS)}
  \label{tab:sidebar_compare}
  \begin{tabularx}{\textwidth}{L{3.5cm}L{4cm}X}
    \toprule
    \thead{Aspect} & \thead{Version 1} & \thead{Version 2 (Sprint 6)} \\
    \midrule
    \rowcolor{lightgray}
    Structure       & Liste plate         & 3 sections : Principal / Collaboration / Administration \\
    Animations      & Aucune              & Framer Motion : collapse/expand, active indicator, label fade \\
    \rowcolor{lightgray}
    Badges          & Aucun               & Notifications (rouge, non-lues) + Audit (vert, nouvelles activités) \\
    Identification  & Aucune              & User card : avatar + nom + rôle + logout \\
    \rowcolor{lightgray}
    Couleur fond    & Variable            & \texttt{bg-slate-950} + bordure \texttt{rgba(255,255,255,0.06)} \\
    \bottomrule
  \end{tabularx}
\end{table}

\subsubsection{Optimisations de performance — Bilan Sprint 6}

\begin{table}[H]
  \centering
  \caption{Optimisations techniques réalisées au cours du Sprint 6}
  \label{tab:perf_s6}
  \begin{tabularx}{\textwidth}{L{4.5cm}L{3cm}X}
    \toprule
    \thead{Problème} & \thead{Impact avant} & \thead{Solution} \\
    \midrule
    \rowcolor{lightgray}
    N+1 GanttController & 1 requête / tâche en cascade &
      Pré-chargement de tous les tasks + dependents en 2 requêtes fixes \\
    Stats audit calculées à chaque page &
      +4 requêtes SQL par appel &
      Paramètre \texttt{?with\_stats=1} optionnel — le polling ne les déclenche jamais \\
    \rowcolor{lightgray}
    AiController::chat — 3 queries projets &
      3 requêtes séparées &
      Scope Eloquent \texttt{accessibleToUser(\$user)} → 1 seule requête \\
    Listener WebSocket manquant &
      Suppressions de commentaires non synchronisées &
      \texttt{.task.comment.deleted} ajouté dans ProjectDetailPage \\
    \rowcolor{lightgray}
    DashboardPage — couverture WebSocket incomplète &
      Données stales après sprint/progression &
      \texttt{.sprint.state.changed} + \texttt{.progress.updated} ajoutés \\
    \bottomrule
  \end{tabularx}
\end{table}

% ============================================================
\subsection{Tests et validation}

\begin{lstlisting}[style=phpstyle,caption={Tests unitaires — RBAC et journal d'audit},
  label={lst:tests_s6}]
// Test : seul un admin peut acceder aux roles
it('denies access to roles for non-admin users', function () {
    $member = User::factory()->create()->assignRole('member');
    $this->actingAs($member)
         ->getJson('/api/roles')
         ->assertStatus(403);
});

// Test : le journal trace les valeurs avant/apres
it('logs old and new values on task update', function () {
    $task = Task::factory()->create(['priority' => 'low']);
    $admin = User::factory()->create()->assignRole('admin');

    $this->actingAs($admin)
         ->patchJson("/api/projects/{$task->project_id}/tasks/{$task->id}",
                     ['priority' => 'high'])
         ->assertOk();

    $this->assertDatabaseHas('activity_logs', [
        'action' => 'task.updated',
    ]);
    $log = ActivityLog::latest()->first();
    expect($log->old_values['priority'])->toBe('low')
     ->and($log->new_values['priority'])->toBe('high');
});

// Test : severity auto-detectee
it('marks deleted task logs as critical', function () {
    $task = Task::factory()->create();
    $admin = User::factory()->create()->assignRole('admin');

    $this->actingAs($admin)
         ->deleteJson("/api/projects/{$task->project_id}/tasks/{$task->id}")
         ->assertOk();

    $this->assertDatabaseHas('activity_logs', [
        'action'   => 'task.deleted',
        'severity' => 'critical',
    ]);
});
\end{lstlisting}

\begin{chapconclusion}
Le sprint~6 élève la plateforme TAC-TIC IntelliPM au niveau enterprise sur deux
dimensions complémentaires. Le système RBAC à permissions atomiques offre une
granularité de contrôle d'accès suffisante pour des organisations multi-équipes
aux exigences variables. Le journal d'audit enrichi — avec ses valeurs avant/après,
sa détection automatique de sévérité, son dashboard analytique aux graphiques
CSS/SVG et son export PDF — constitue un outil de gouvernance à part entière.
Les optimisations de performance réalisées en parallèle (élimination du N+1 Gantt,
stats optionnelles, couverture WebSocket complète) garantissent la fiabilité de
l'architecture en conditions de charge réelle. Avec ce sprint, la plateforme
atteint un niveau de maturité production-ready sur l'ensemble de ses 34 user
stories et 219 story points livrés.
\end{chapconclusion}
